\section{Waiver and Repair Templates}

Waivers are allowed only when they are named, scoped, owned, testable, and temporary enough for an agent to repair later. Runtime exceptions are program errors; standards deviations are waivers.

\begin{lstlisting}[style=codebox]
---
id: WVR-YYYY-NNN
state: proposed | active | renewed | expired | closed | superseded
owner: <team or owner cell>
scope:
  - <exact file, crate, service, contract, or zone>
standard_version: "0.8.0"
auditor_version: "0.8.8"
schema_version: "1.5.0"
reason: <why the normal standard cannot be followed now>
compensating_control: <alternate evidence or manual gate>
required_proof:
  - <test-map lane name and command>
residual_risk: <bounded risk statement>
expires_at: <date or measurable exit condition>
owner_approval: <receipt or reviewer>
---

# WVR-YYYY-NNN: <short name>

## Repair path
- <first bounded repair pattern>
- <second bounded repair pattern>
\end{lstlisting}

\begin{lstlisting}[style=codebox]
---
receipt_id: RPR-YYYY-NNN
standard_version: "0.8.0"
auditor_version: "0.8.8"
schema_version: "1.5.0"
before_commit: <sha>
after_commit: <sha>
started_at: <timestamp>
completed_at: <timestamp>
command_digest: sha256:...
log_digest: sha256:...
artifact_digests:
  agent/repo-score.json: sha256:...
decision: pass | review | block | ratchet_fail | release_fail
waiver_closed: <id or false>
residual_risk: <none or bounded risk>
redaction_status: clean | redacted | restricted
owner_approval: <receipt or reviewer>
expires_at: <required if partial>
---

# Repair receipt

Changed paths, proof lane, result, and next repair.
\end{lstlisting}
